% Authority: exact-scope leaf 6; full PDF page 85; printed p.68.
% U:p068.u01 | continuation of p067.u03
\noindent
of the form $4n+1$, $\leg{k}{p}=1$. One finds in the same way:
\[
\leg{k'}{p}=1,\quad \leg{k''}{p}=1,\quad\text{etc.}
\]
On the other hand, since $p$ is of the form $8n+1$, one also has $\leg{2}{p}=1$,
and consequently $\leg{2^\nu}{p}=1$. Multiplying this last relation by all the
preceding ones, one obtains:
\[
\leg{2^\nu kk'k''\ldots}{p}=\leg{u}{p}=1.
\]
% U:p068.u02 | continuation
Now consider the simple factors of the odd number $t$, which we shall divide
into two classes. The first class will comprise the prime divisors of one of
these two forms: $8n+1$, $8n+7$, and the numbers belonging to it will be designated
by $g$, $g'$, $g''$, $\ldots$; the second class will be composed of numbers
$h$, $h'$, $h''$, $\ldots$, contained in these two forms: $8n+3$, $8n+5$.
One first has:
\[
t=gg'g''\ldots\times hh'h''\ldots,
\]
and one will then conclude from the equation $p=t^2+2u^2$:
\[
\begin{aligned}
\leg{2p}{g}&=1, &\leg{2p}{g'}&=1, &\leg{2p}{g''}&=1, &\text{etc.}\\
\leg{2p}{h}&=1, &\leg{2p}{h'}&=1, &\leg{2p}{h''}&=1, &\text{etc.}
\end{aligned}
\]
% U:p068.u03 | continuation
On the other hand, by known theorems one has:
\[
\begin{aligned}
\leg{2}{g}&=1, &\leg{2}{g'}&=1, &\leg{2}{g''}&=1, &\text{etc.}\\
\leg{2}{h}&=-1, &\leg{2}{h'}&=-1, &\leg{2}{h''}&=-1, &\text{etc.}
\end{aligned}
\]
% U:p068.u04 | continuation
If one now compares these relations with the preceding ones, one will find:
\[
\begin{aligned}
\leg{p}{g}&=1, &\leg{p}{g'}&=1, &\leg{p}{g''}&=1, &\text{etc.}\\
\leg{p}{h}&=-1, &\leg{p}{h'}&=-1, &\leg{p}{h''}&=-1, &\text{etc.}
\end{aligned}
\]
% U:p068.u05 | continuation
The application of the law of reciprocity to these last relations gives:
\[
\begin{aligned}
\leg{g}{p}&=1, &\leg{g'}{p}&=1, &\leg{g''}{p}&=1, &\text{etc.}\\
\leg{h}{p}&=-1, &\leg{h'}{p}&=-1, &\leg{h''}{p}&=-1, &\text{etc.}
\end{aligned}
\]
